\documentclass[11pt,a4paper]{article}

% ==========================================
% Packages & Formatting
% ==========================================
\usepackage[utf8]{inputenc}
\usepackage[T1]{fontenc}
\usepackage{amsmath, amssymb, amsfonts, bm}
\usepackage{geometry}
\usepackage{booktabs}
\usepackage{enumitem}
\usepackage{siunitx}
\usepackage{comment}
\geometry{margin=2.5cm}

\usepackage{xcolor} % <-- Agrega esta línea
\usepackage{hyperref}
\hypersetup{
    colorlinks=true,
    linkcolor=blue!70!black,
    urlcolor=blue!70!black,
    citecolor=blue!70!black
}

% Custom Vector/Tensor Macros
\newcommand{\vect}[1]{\bm{#1}}
\newcommand{\dd}{\mathrm{d}}
\newcommand{\pderiv}[2]{\frac{\partial #1}{\partial #2}}
\newcommand{\mderiv}[2]{\frac{\mathrm{D} #1}{\mathrm{D} #2}}
\newcommand{\Div}{\boldsymbol{\nabla}\!\cdot}
\newcommand{\Curl}{\boldsymbol{\nabla}\!\times}
\newcommand{\grad}{\boldsymbol{\nabla}}
\newcommand{\uvec}{\boldsymbol{u}}
\newcommand{\Bvec}{\boldsymbol{B}}
\newcommand{\Evec}{\boldsymbol{E}}
\newcommand{\Jvec}{\boldsymbol{J}}
\newcommand{\Svec}{\boldsymbol{S}}
\newcommand{\rhov}{\rho}
\newcommand{\ppl}{p}

\title{\textbf{Multi-Scale Plasma Physics of Solar Wind Interactions: From Kinetic-MHD Foundations to Reduced Engineering and Astrophysical Models}}
\author{Space Plasma Physics \& Spacecraft Systems Documentation}
\date{\today}

\begin{document}

\maketitle

\begin{abstract}
This document establishes a unified theoretical framework governing the interaction of the solar wind plasma with artificial spacecraft systems and natural cometary bodies. Starting from the overarching kinetic-fluid hierarchy—specifically the Vlasov--Boltzmann--Maxwell system and Multi-Species Generalized Magnetohydrodynamics (MHD)—we systematically derive reduced-order mathematical models tailored to three distinct regimes: (1) active electromagnetic radiation shielding via magnetostatics and test-particle Lorentz dynamics, and (2) solar wind--comet plasma tail interactions incorporating mass loading, pickup ions, and Alfvénic draping.
\end{abstract}

\tableofcontents
\newpage

% ==========================================
% Nomenclature
% ==========================================
\section*{Nomenclature}
\addcontentsline{toc}{section}{Nomenclature}

\begin{center}
\begin{tabular}{lp{5cm} @{\hspace{1.2cm}} lp{5cm}}
\toprule
\textbf{Symbol} & \textbf{Description} & \textbf{Symbol} & \textbf{Description} \\
\midrule
$\epsilon_0$ & Vacuum permittivity & $\mu_0$ & Vacuum permeability \\
$c$ & Speed of light in vacuum & $k_B$ & Boltzmann constant \\
$s$ & Particle species subscript ($e, p, i$) & $t$ & Time \\
$q_s, Z_s e$ & Species electric charge & $m_s$ & Species particle mass \\
$n_s$ & Particle number density & $\rho_s = m_s n_s$ & Mass density \\
$\vect{v}_s, \vect{u}_s$ & Microscopic / bulk fluid velocity & $\vect{p}_s$ & Relativistic momentum \\
$\vect{E}$ & Electric field vector & $\vect{B}$ & Magnetic flux density vector \\
$\vect{J}$ & Electric current density & $\eta, \sigma$ & Electrical resistivity / conductivity \\
$\mathbf{P}_s$ & Pressure tensor ($P_\parallel, P_\perp$) & $\boldsymbol{\sigma}$ & Total viscous/electromagnetic stress tensor \\
$\psi_s(\vect{r},\vect{v},t)$ & Phase-space distribution function & $\Phi$ & Viscous dissipation function \\
$S_{m,s}$ & Mass source/sink term (photoionization) & $\vect{S}_{p,s}$ & Inter-species collisional momentum transfer \\
$S_{e,s}$ & Volumetric heating source & $\vect{Q}_s$ & Heat flux vector \\
$\lambda_D$ & Debye shielding length & $r_{L,s}$ & Larmor (gyration) radius \\
$d_i$ & Ion inertial length ($c/\omega_{pi}$) & $v_A$ & Alfvén speed \\
$\beta$ & Plasma thermal-to-magnetic pressure ratio & $\gamma$ & Lorentz factor / adiabatic index \\
\bottomrule
\end{tabular}
\end{center}

% ==========================================
% SECTION 1: UNIFIED THEORETICAL FRAMEWORK
% ==========================================
\section{Comprehensive Framework: Kinetic-MHD Foundations of the Solar Wind}

The solar wind is a weakly collisional, magnetized, quasi-neutral plasma streaming radially outward from the solar corona at supersonic and super-Alfvénic speeds ($v_{\text{sw}} \sim 300\text{--}800\;\si{km/s}$ at $1\;\si{AU}$, reaching $>1000\;\si{km/s}$ during fast coronal mass ejections). Capturing all physical phenomena across spatial scales—from sub-millimeter electrostatic sheaths to astronomical-unit cometary magnetotails—requires starting from the fundamental kinetic-electromagnetic equations and systematically taking velocity moments to obtain Generalized Magnetohydrodynamics (MHD).

\subsection{Maxwell's Equations}
The self-consistent electromagnetic fields $\vect{E}(\vect{r},t)$ and $\vect{B}(\vect{r},t)$ generated by free charges and currents in the plasma are governed by Maxwell's equations:
\begin{align}
    \nabla \cdot \vect{E} &= \frac{1}{\epsilon_0} \sum_s q_s n_s, \label{eq:gauss_e} \\
    \nabla \cdot \vect{B} &= 0, \label{eq:gauss_b} \\
    \nabla \times \vect{E} &= -\pderiv{\vect{B}}{t}, \label{eq:faraday} \\
    \nabla \times \vect{B} &= \mu_0 \left( \vect{J} + \epsilon_0 \pderiv{\vect{E}}{t} \right), \quad \text{where } \vect{J} = \sum_s q_s n_s \vect{u}_s. \label{eq:ampere}
\end{align}

\subsection{Kinetic Description: Vlasov--Boltzmann--Fokker--Planck Equation}
At microscopic scales comparable to the Debye length ($\lambda_D$) or ion gyroradius ($r_{L,i}$), each species $s$ is described by a single-particle phase-space distribution function $\psi_s(\vect{r}, \vect{p}_s, t)$. Its evolution is governed by the relativistic Boltzmann equation with a Fokker--Planck collision operator for long-range Coulomb interactions:
\begin{equation}
    \pderiv{\psi_s}{t} + \vect{v}_s \cdot \nabla \psi_s + q_s \left( \vect{E} + \vect{v}_s \times \vect{B} \right) \cdot \nabla_{\vect{p}_s} \psi_s = -\pderiv{}{v_{s,i}} \left( \psi_s \langle \Delta v_{s,i} \rangle \right) + \frac{1}{2} \frac{\partial^2}{\partial v_{s,i} \partial v_{s,j}} \left( \psi_s \langle \Delta v_{s,i} \Delta v_{s,j} \rangle \right) + \left( \pderiv{\psi_s}{t} \right)_{\text{source}},
    \label{eq:boltzmann_fp}
\end{equation}
where $\vect{v}_s = \vect{p}_s / (\gamma_s m_s)$ and $\gamma_s = \left(1 + \|\vect{p}_s\|^2 / (m_s^2 c^2)\right)^{1/2}$. When collisions are negligible ($\nu_{ei} \tau_{\text{dyn}} \ll 1$, valid for the ambient solar wind where the mean free path $\lambda_{\text{mfp}} \sim 1\;\si{AU}$), the right-hand side vanishes, yielding the collisionless \textbf{Vlasov equation}.

\subsection{Multi-Species Generalized MHD (Moment Equations)}
Integrating Eq.~\eqref{eq:boltzmann_fp} over velocity space ($\int \dd^3 v$, $\int m_s \vect{v} \dd^3 v$, and $\int \frac{1}{2} m_s v^2 \dd^3 v$) produces the macroscopic fluid conservation laws for mass, momentum, and energy for each species $s$:

\begin{align}
    \pderiv{\rho_s}{t} + \nabla \cdot (\rho_s \vect{u}_s) &= S_{m,s}, \label{eq:gen_mass} \\
    \pderiv{(\rho_s \vect{u}_s)}{t} + \nabla \cdot \left( \rho_s \vect{u}_s \vect{u}_s + \mathbf{P}_s \right) &= q_s n_s \left( \vect{E} + \vect{u}_s \times \vect{B} \right) + \vect{S}_{p,s} + S_{m,s} \vect{u}_{s,0}, \label{eq:gen_mom} \\
    \pderiv{}{t} \left( \frac{1}{2} \rho_s u_s^2 + \mathcal{E}_s \right) + \nabla \cdot \left[ \left( \frac{1}{2} \rho_s u_s^2 + \mathcal{E}_s \right) \vect{u}_s + \mathbf{P}_s \cdot \vect{u}_s + \vect{Q}_s \right] &= q_s n_s \vect{u}_s \cdot \vect{E} + S_{e,s} \label{eq:gen_energy}
\end{align}
where $\rho_s = m_s n_s$, $\mathcal{E}_s$ is the internal thermal energy density, and $\mathbf{P}_s$ is the gyrotropic anisotropic pressure tensor aligned with the local magnetic field unit vector $\hat{\vect{b}} = \vect{B}/\|\vect{B}\|$:
\begin{equation}
    \mathbf{P}_s = P_{\perp,s} \mathbf{I} + \left( P_{\parallel,s} - P_{\perp,s} \right) \hat{\vect{b}}\hat{\vect{b}}^T.
    \label{eq:anisotropic_p}
\end{equation}

\subsection{Single-Fluid Closure and Generalized Ohm's Law}
Summing Eqs.~\eqref{eq:gen_mass}--\eqref{eq:gen_energy} over electrons and ions under quasi-neutrality ($n_e \approx \sum_i Z_i n_i$) and neglecting displacement currents ($\epsilon_0 \partial_t \vect{E} \ll \vect{J}$) yields single-fluid resistive/Hall MHD. Subtracting the electron momentum equation from the ion momentum equation (assuming $m_e \ll m_i$) yields the **Generalized Ohm's Law**:
\begin{equation}
    \vect{E} + \vect{u} \times \vect{B} = \eta \vect{J} + \frac{1}{n_e e} \left( \vect{J} \times \vect{B} - \nabla \cdot \mathbf{P}_e \right) + \frac{m_e}{n_e e^2} \left[ \pderiv{\vect{J}}{t} + \nabla \cdot \left( \vect{u}\vect{J} + \vect{J}\vect{u} \right) \right].
    \label{eq:gen_ohm}
\end{equation}
Here, the right-hand side separates non-ideal plasma effects by spatial scale:
\begin{enumerate}
    \item $\eta \vect{J}$: Ohmic resistivity (diffusive magnetic reconnection).
    \item $\frac{1}{n_e e} \vect{J} \times \vect{B}$: Hall term, dominant at scales comparable to the ion inertial length $d_i = c/\omega_{pi}$.
    \item $-\frac{1}{n_e e} \nabla \cdot \mathbf{P}_e$: Ambipolar electron pressure gradient, governing electrostatic potential barriers and double layers.
    \item $\frac{m_e}{n_e e^2} \partial_t \vect{J}$: Electron inertia, active at the electron skin depth $d_e = c/\omega_{pe}$.
\end{enumerate}

When all non-ideal terms vanish ($\vect{E} + \vect{u} \times \vect{B} = \vect{0}$), magnetic flux is frozen into the plasma fluid (**Ideal MHD**). Together with Poynting's theorem,
\begin{equation}
    \pderiv{}{t} \left( \frac{\epsilon_0 E^2}{2} + \frac{B^2}{2\mu_0} \right) + \nabla \cdot \left( \frac{\vect{E} \times \vect{B}}{\mu_0} \right) = -\vect{J} \cdot \vect{E},
    \label{eq:poynting}
\end{equation}
this system describes the global energy exchange between electromagnetic fields and bulk plasma kinetic/thermal energy.

\subsection{Characteristic Plasma Dimensionless Parameters}
To transition rigorously from Eqs.~\eqref{eq:gauss_e}--\eqref{eq:gen_ohm} to engineering and astrophysical sub-models, three key dimensionless parameters dictate which terms may be neglected:
\begin{itemize}
    \item \textbf{Plasma Beta ($\beta$):} Ratio of thermal pressure to magnetic pressure, $\beta = \frac{2\mu_0 n k_B T}{B^2}$. In the ambient solar wind, $\beta \sim 0.5\text{--}1$; near superconducting spacecraft coils, $\beta \ll 1$ (magnetically dominated).
    \item \textbf{Magnetic Reynolds Number ($R_m$):} Ratio of magnetic advection to resistive diffusion, $R_m = \frac{\mu_0 u L}{\eta}$. For large-scale cometary tails ($L \sim 10^6\;\si{km}$), $R_m \gg 1$ (frozen-in ideal MHD holds).
    \item \textbf{Ion Magnetization / Hall Parameter ($\delta_i = d_i / L$):} Ratio of ion inertial length (or gyroradius $r_{L,i}$) to the characteristic obstacle scale $L$. For a $1\;\si{U}$ CubeSat ($L = 0.1\;\si{m}$), $r_{L,i} \gg L$, meaning ions are unmagnetized on the spacecraft scale and must be treated kinetically rather than as a continuum MHD fluid.
\end{itemize}

---

% ==========================================
% SECTION 2: ELECTROMAGNETIC SHIELDS
% ==========================================
\section{Active Electromagnetic Shields for Spacecraft}

\subsection{Scale-Reduction Assumptions from Kinetic-MHD}
When designing an active magnetic shield for a small spacecraft (such as a $1\;\si{U}$ CubeSat with $L \approx 0.1\;\si{m}$) using high-temperature superconducting (YBCO) coils, the characteristic system scale $L$ is orders of magnitude smaller than the solar wind ion inertial length ($d_i \approx 100\;\si{km}$) and Debye length ($\lambda_D \approx 10\;\si{m}$):
\begin{equation}
    L \ll \lambda_D \ll d_i \sim r_{L,i}.
\end{equation}
Consequently, the single-fluid MHD continuum approximation breaks down ($\delta_i \gg 1$). The comprehensive system reduces via the following hierarchy of approximations:
\begin{enumerate}
    \item \textbf{Negligible Plasma Back-Reaction (Vacuum Field Limit):} Because the ambient solar wind current density is tiny compared to the superconducting coil current density ($J_{\text{sw}} \ll J_{\text{coil}}$) and $\beta_{\text{local}} \ll 1$ inside the standoff cavity, the induced plasma magnetic field is neglected ($\vect{B} \approx \vect{B}_{\text{coil}} + \vect{B}_{\text{IMF}}$).
    \item \textbf{Magnetostatic Regime:} Coil currents are steady ($\partial_t \vect{J}_{\text{coil}} = 0$), reducing Ampère's law (Eq.~\ref{eq:ampere}) to $\nabla \times \vect{B} = \mu_0 \vect{J}_{\text{coil}}$.
    \item \textbf{Test-Particle Kinetic Limit:} Rather than solving the Vlasov equation for the entire distribution $\psi_s$, the characteristics (trajectories) of individual high-energy solar wind protons and electrons are integrated directly using the Lorentz force equation.
\end{enumerate}

\subsection{Magnetostatic Field Formulation (Biot--Savart Law)}
Under the magnetostatic condition $\nabla \times \vect{B} = \mu_0 \vect{J}$ and $\nabla \cdot \vect{B} = 0$, introducing the magnetic vector potential $\vect{B} = \nabla \times \vect{A}$ in the Coulomb gauge ($\nabla \cdot \vect{A} = 0$) yields Poisson's equation $\nabla^2 \vect{A} = -\mu_0 \vect{J}$. Its Green's function solution gives the **Biot--Savart Law** for a volumetric conductor of current density $\vect{J}(\vect{r}')$:
\begin{equation}
    \vect{B}(\vect{r}) = \frac{\mu_0}{4\pi} \iiint_{V'} \frac{\vect{J}(\vect{r}') \times (\vect{r} - \vect{r}')}{\|\vect{r} - \vect{r}'\|^3} \, \dd^3 r'.
    \label{eq:biot_savart_3d}
\end{equation}
In spectral (wave-number) space $\vect{k}$, this convolution simplifies algebraically to:
\begin{equation}
    \hat{\vect{B}}(\vect{k}) = \frac{i \mu_0}{\|\vect{k}\|^2} \left( \vect{k} \times \hat{\vect{J}}(\vect{k}) \right),
\end{equation}
allowing rapid 3D field evaluation over multi-coil configurations (such as a 4-arm, 3-coil-per-arm pumpkin topology). Within the superconducting YBCO windings, a small resistive flux-flow electric field arises near the critical current density $J_c$, modeled via the empirical $E\text{--}J$ power law:
\begin{equation}
    \vect{E}_{\text{coil}}(\vect{r}) = E_c \left( \frac{\|\vect{J}(\vect{r})\|}{J_c} \right)^{n-1} \vect{J}(\vect{r}),
\end{equation}
where $E_c$ is the critical electric field criterion ($1\;\si{\mu V/cm}$) and $n$ is the flux-creep exponent.

\subsection{Relativistic Lorentz Force Trajectory Dynamics}
Each incident solar wind ion (charge $q$, rest mass $m$) is propagated as a Lagrangian test particle along the characteristics of the collisionless Vlasov equation:
\begin{align}
    \frac{\dd \vect{r}}{\dd t} &= \vect{v}, \\
    \frac{\dd (\gamma m \vect{v})}{\dd t} &= q \left( \vect{E}_{\text{coil}}(\vect{r}) + \vect{v} \times \vect{B}_{\text{coil}}(\vect{r}) \right), \quad \gamma = \left(1 - \frac{\|\vect{v}\|^2}{c^2}\right)^{-1/2}.
    \label{eq:lorentz_ode}
\end{align}
For typical solar wind protons ($v_{\text{sw}} \approx 400\text{--}1000\;\si{km/s}$), $\gamma \approx 1.0000056$, allowing the non-relativistic simplification $\dd \vect{v}/\dd t = (q/m)(\vect{E} + \vect{v} \times \vect{B})$, whereas relativistic $\gamma$ must be retained for solar energetic particles (SEPs) in the $\si{MeV}$--$\si{GeV}$ range.

\subsection{Magnetic Standoff Criterion}
To prevent solar wind protons from reaching the protected spacecraft core of characteristic radius $R_{\text{prot}}$, the magnetic pressure of the shield must balance the dynamic ram pressure of the incoming solar wind (derived from integrating the momentum conservation equation Eq.~\ref{eq:gen_mom} across the magnetopause boundary):
\begin{equation}
    \frac{B_{\text{eq}}^2}{2\mu_0} \approx \rho_{\text{sw}} v_{\text{sw}}^2 = m_p n_{\text{sw}} v_{\text{sw}}^2 \implies B_{\text{eq}} = \sqrt{2 \mu_0 m_p n_{\text{sw}} v_{\text{sw}}^2}.
\end{equation}
Additionally, the minimum integrated magnetic field impulse $\int B_\perp \dd l$ across the shielding layer must exceed the proton magnetic rigidity $R_m = \gamma m_p v_{\text{sw}} / q_p$ so that the local Larmor radius $r_{L,p} = \frac{\gamma m_p v_{\text{sw}}}{q_p B}$ is smaller than the physical thickness of the magnetic barrier.

---

% ==========================================
% SECTION 3: ELECTRIC SOLAR WIND SAILS
% ==========================================
\begin{comment}
\section{Electric Solar Wind Sails (E-Sail)}

\subsection{Scale-Reduction Assumptions from Kinetic-MHD}
An Electric Solar Wind Sail (E-Sail) harvests momentum from the solar wind using long, thin, positively biased conducting tethers (wire radius $r_w \sim 25\text{--}50\;\si{\mu m}$, length $L_t \sim 1\text{--}20\;\si{km}$, potential $V_0 \sim +10\text{--}20\;\si{kV}$). Comparing physical scales:
\begin{equation}
    r_w \ll \lambda_D \ll L_t \ll d_i.
\end{equation}
Because the tether diameter is much smaller than the Debye length ($\lambda_D \sim 10\;\si{m}$), an electrostatic **electron-depleted Debye sheath** forms around each tether with an effective sheath radius $r_s \sim \text{several } \lambda_D$. The full Vlasov--Maxwell system reduces via:
\begin{enumerate}
    \item \textbf{Electrostatic Dominance:} The tether current $I_t$ is small enough that its self-generated magnetic field $B_t = \mu_0 I_t / (2\pi r)$ exerts a negligible Lorentz force compared to the radial electrostatic field $\vect{E} = -\nabla \phi$ ($\|\vect{v} \times \vect{B}\| \ll \|\vect{E}\|$).
    \item \textbf{Vlasov--Poisson Sheath Equilibrium:} Electrons are trapped in the positive potential well while incoming solar wind protons are elastically reflected/deflected by the potential barrier $\phi(r)$ governed by Poisson's equation (Eq.~\ref{eq:gauss_e}), $\nabla^2 \phi = -\frac{e}{\epsilon_0}(n_i - n_e)$.
    \item \textbf{Macroscopic Momentum-Energy Conservation:} Rather than resolving every ion trajectory in 3D, thrust is computed via global conservation of linear momentum and kinetic energy across a control volume enclosing the tether sheath, integrated over the solar wind proton flux $n_p v_{\text{sw}}$.
\end{enumerate}

\subsection{Elastic Collision \& Impulse-Momentum Derivation of $\Delta V$}
Consider a single solar wind proton of mass $m_p$ and initial velocity $v_{p,i}$ interacting elastically with a spacecraft/sail of mass $m_v$ and initial velocity $v_{v,i}$ (in the heliocentric radial frame). Conservation of linear momentum and total mechanical energy yield:
\begin{align}
    m_p v_{p,i} + m_v v_{v,i} &= m_p v_{p,f} + m_v v_{v,f}, \label{eq:esail_mom} \\
    \frac{1}{2} m_p v_{p,i}^2 + \frac{1}{2} m_v v_{v,i}^2 &= \frac{1}{2} m_p v_{p,f}^2 + \frac{1}{2} m_v v_{v,f}^2. \label{eq:esail_en}
\end{align}
Eliminating the post-interaction proton velocity $v_{p,f}$ between Eqs.~\eqref{eq:esail_mom} and \eqref{eq:esail_en} yields the final sail velocity $v_{v,f}$:
\begin{equation}
    v_{v,f} = \frac{2 m_p v_{p,i} + (m_v - m_p) v_{v,i}}{m_v + m_p}.
\end{equation}
Thus, the incremental velocity change ($\delta V = v_{v,f} - v_{v,i}$) imparted by a single head-on proton reflection is:
\begin{equation}
    \delta V = \frac{2 (v_{p,i} - v_{v,i})}{1 + \frac{m_v}{m_p}}.
    \label{eq:single_dv}
\end{equation}
*(Note: When $v_{p,i}$ and $v_{v,i}$ are defined with opposite sign conventions for counter-streaming frames, the numerator takes the form $2(v_{p,i} + v_{v,i})$).*

More generally, when accounting for an electromagnetic impulse $I_{\text{EM}} = \int \vect{F} \, \dd t$ and electromagnetic field work $E_{\text{EM}} = \int \vect{F} \cdot \dd \vect{r}$ delivered over a finite interaction time $\Delta t$, the momentum and energy balances generalize to $p = m_p v_{p,i} + m_v v_{v,i} + I_{\text{EM}}$ and $e = m_v v_{v,i}^2 + m_p v_{p,i}^2 + 2E_{\text{EM}}$, yielding:
\begin{equation}
    v_{v,f} = \frac{p \pm \sqrt{p^2 + \left(1 + \frac{m_p}{m_v}\right)(m_p e - p^2)}}{m_p + m_v}.
\end{equation}

\subsection{Continuous Thrust, Force per Unit Length, and Total $\Delta V$}
In a continuous solar wind stream of proton number density $n_p$ and relative flow speed $v_{\text{rel}} = v_{\text{sw}} - v_v$, the rate of protons intercepted by an effective sheath cross-sectional area $A_{\text{eff}} = 2 r_s L_{\text{total}}$ (for $N$ tethers of length $L$, $L_{\text{total}} = NL$) is $\dot{N}_p = n_p v_{\text{rel}} A_{\text{eff}}$. Multiplying Eq.~\eqref{eq:single_dv} by $\dot{N}_p$ in the limit $m_v \gg m_p$ yields the continuous thrust force:
\begin{equation}
    F = m_v \frac{\dd v_v}{\dd t} = 2 m_p n_p v_{\text{rel}}^2 A_{\text{eff}} \eta_{\text{defl}},
\end{equation}
where $\eta_{\text{defl}}$ is the deflection efficiency factor.

From kinetic-Vlasov sheath theory (Janhunen's E-Sail formulation), the electrostatic force per unit tether length $\sigma_f$ scales with the tether potential $V_0$ (corrected for solar wind motional potential $V_w$) and the Debye-length sheath width:
\begin{equation}
    \sigma_f = K \, m_p n_p v_{\text{rel}}^2 \, r_s \approx 0.18 \, (V_0 - V_w) v_{\text{rel}} \sqrt{\epsilon_0 m_p n_p},
    \label{eq:esail_force_per_length}
\end{equation}
where $K \approx 3.09$ is a dimensionless sheath geometry constant. Because the solar wind density decreases with heliocentric distance $r$ as $n_p(r) = n_\oplus (r_\oplus / r)^2$ (from steady spherical mass conservation $\nabla \cdot (n \vect{v}) = 0$ at constant $v_{\text{sw}}$), the force per unit length scales as $\sigma_f(r) \propto \sqrt{n_p(r)} \propto r^{-1}$ (or $r^{-7/6}$ when accounting for radial electron temperature cooling $T_e(r) \propto r^{-2/3}$ in $\lambda_D$).

Consequently, the characteristic acceleration $a_c$ at $1\;\si{AU}$ ($r_\oplus$) and the radial heliocentric acceleration vector $\vect{a}(r)$ are:
\begin{equation}
    a_c = \frac{N L \, \sigma_{f,\oplus}}{m_v}, \qquad \vect{a}(r) = \frac{a_c \tau}{2} \left( \frac{r_\oplus}{r} \right) \left[ \hat{\vect{r}} + (\hat{\vect{r}} \cdot \hat{\vect{n}})\hat{\vect{n}} \right],
    \label{eq:esail_accel}
\end{equation}
where $\tau \in [0,1]$ is the voltage throttling factor and $\hat{\vect{n}}$ is the unit normal to the sail spin plane. Integrating the tangential and radial acceleration over mission duration $T$ gives the total mission $\Delta V$:
\begin{equation}
    \Delta V = \int_0^T \|\vect{a}(t)\| \, \dd t \approx \sum_{k} \frac{N L \, \sigma_{f,k}}{m_v} \Delta t_k.
\end{equation}
\end{comment}

---

% ==========================================
% SECTION 4: SOLAR WIND - COMET INTERACTION
% ==========================================
\section{Solar Wind -- Comet Plasma Tail Interaction}

\subsection{Scale-Reduction Assumptions from Kinetic-MHD}
Unlike a spacecraft, an active comet outgasses neutral volatiles (predominantly $\text{H}_2\text{O}$, $\text{CO}$, and $\text{CO}_2$) at a production rate $Q_n \sim 10^{28}\text{--}10^{30}\;\si{molecules/s}$, forming an extended neutral coma spanning $10^5\text{--}10^6\;\si{km}$. Here, the interaction scale $L_{\text{comet}}$ greatly exceeds both the Debye length and the solar wind proton inertial length:
\begin{equation}
    \lambda_D \ll d_i \ll L_{\text{comet}}, \qquad R_m \gg 1.
\end{equation}
Because $R_m \gg 1$ and $L_{\text{comet}} \gg d_i$, the full system reduces to **Single-Fluid Ideal MHD with Source Terms (Mass-Loaded Ideal MHD)**, coupled to a neutral gas continuity model:
\begin{enumerate}
    \item \textbf{Frozen-In Magnetic Field:} Resistive and Hall terms in Generalized Ohm's Law (Eq.~\ref{eq:gen_ohm}) are neglected outside thin current sheets, giving $\vect{E} = -\vect{u} \times \vect{B}$.
    \item \textbf{Photoionization \& Charge-Exchange Mass Loading:} Neutral cometary molecules expand radially uninhibited by electromagnetic fields until solar EUV photons or charge-exchange collisions ionize them (e.g., $\text{H}_2\text{O} + h\nu \to \text{H}_2\text{O}^+ + e^-$, $\text{CO} + h\nu \to \text{CO}^+ + e^-$), injecting heavy ions at nearly zero velocity into the fast solar wind flow ($S_m > 0$).
    \item \textbf{Diffusion Approximation in the Ion Tail:} Far downstream in the Type I plasma tail, plasma transport perpendicular and parallel to the draped magnetic field is governed by advection-diffusion equations (such as the Shulman--Nazarchuk diffusion model).
\end{enumerate}

\begin{comment}
\subsection{Neutral Coma Outgassing (Haser Model)}
Before ionization, neutral molecules expand radially from the nucleus at velocity $v_n \approx 0.8\text{--}1.0\;\si{km/s}$ while undergoing photodissociation and photoionization with characteristic destruction timescale $\tau_d$ (scale length $\lambda_n = v_n \tau_d$). Steady spherical mass conservation (Eq.~\ref{eq:gen_mass} with sink $-n_n / \tau_d$) yields the **Haser Neutral Density Profile**:
\begin{equation}
    n_n(r) = \frac{Q_n}{4\pi r^2 v_n} \exp\left( -\frac{r}{\lambda_n} \right).
    \label{eq:haser}
\end{equation}
Radio and optical observations (e.g., $\text{OH}$ $18\;\si{cm}$ transitions or $\text{CO}^+$ fluorescence) directly constrain $Q_n$ and the local volumetric ion production rate $S_m$:
\begin{equation}
    S_m(r) = m_i \nu_{\text{ion}} n_n(r) = \frac{m_i Q_n}{4\pi r^2 \lambda_{\text{ion}}} \exp\left( -\frac{r}{\lambda_n} \right),
\end{equation}
where $m_i$ is the mean cometary ion mass ($\approx 18\text{--}28\;\si{amu}$) and $\nu_{\text{ion}} = 1/\tau_{\text{ion}}$ is the photoionization rate.

\subsection{Mass-Loaded Ideal MHD Conservation Equations}
Incorporating the cometary mass source $S_m$, ion-neutral collisional drag $\vect{S}_p = -\nu_{in} \rho (\vect{u} - \vect{u}_n)$, and radiative/photoelectron heating $S_e$ into the single-fluid MHD equations (Eqs.~\ref{eq:gen_mass}--\ref{eq:gen_energy}) yields the **Mass-Loaded Ideal MHD System**:
\begin{align}
    \pderiv{\rho}{t} + \nabla \cdot (\rho \vect{u}) &= S_m, \label{eq:comet_mass} \\
    \pderiv{(\rho \vect{u})}{t} + \nabla \cdot \left[ \rho \vect{u}\vect{u} + \left( P + \frac{B^2}{2\mu_0} \right)\mathbf{I} - \frac{\vect{B}\vect{B}}{\mu_0} \right] &= S_m \vect{u}_n - \nu_{in} \rho (\vect{u} - \vect{u}_n), \label{eq:comet_mom} \\
    \pderiv{\mathcal{E}_{\text{tot}}}{t} + \nabla \cdot \left[ \left( \mathcal{E}_{\text{tot}} + P + \frac{B^2}{2\mu_0} \right)\vect{u} - \frac{(\vect{u} \cdot \vect{B})\vect{B}}{\mu_0} \right] &= \frac{1}{2} S_m u_n^2 + S_e, \label{eq:comet_energy} \\
    \pderiv{\vect{B}}{t} - \nabla \times (\vect{u} \times \vect{B}) &= \vect{0}, \label{eq:comet_induction}
\end{align}
where the total MHD energy density is $\mathcal{E}_{\text{tot}} = \frac{1}{2}\rho u^2 + \frac{P}{\gamma - 1} + \frac{B^2}{2\mu_0}$.

\subsection{Plasma Deceleration, Bow Shock, and Magnetic Field Draping}
Consider a 1D steady flow along the Sun--comet stagnation line ($x$-axis) neglecting magnetic pressure far upstream. Integrating Eqs.~\eqref{eq:comet_mass} and \eqref{eq:comet_mom} shows that adding stationary mass ($u_n \approx 0$) increases the mass flux $\rho u$ while conserving momentum flux $\rho u^2 + P$:
\begin{equation}
    \frac{\dd (\rho u)}{\dd x} = S_m(x) > 0, \qquad \frac{\dd (\rho u^2 + P)}{\dd x} \approx 0.
\end{equation}
Because mass flux increases while momentum flux remains constant, the bulk plasma velocity $u(x)$ must **decelerate continuously** as it approaches the nucleus. When the normalized mass flux reaches a critical threshold (at local Mach number $M = 1$, corresponding to a $\sim 33\%$ increase in mass flux for $\gamma = 5/3$), a smooth supersonic transition becomes impossible and a weak **cometary bow shock** forms upstream of the nucleus, followed closer to the nucleus by a **diamagnetic cavity** (contact surface) where inward magnetic tension balances outward ion-neutral drag:
\begin{equation}
    \frac{B^2}{\mu_0 R_c} \approx \nu_{in} \rho u_n \implies R_{\text{cavity}} \propto \frac{Q_n^{3/4}}{B_{\text{IMF}}}.
\end{equation}

Simultaneously, the induction equation (Eq.~\ref{eq:comet_induction}) dictates that Interplanetary Magnetic Field (IMF) lines frozen into the decelerating central plasma lag behind the outer unperturbed solar wind flow. This differential advection wraps the IMF lines around the cometary ionosphere like a hairpin (**Alfvén magnetic draping**), creating two antiparallel magnetic lobes separated by a central neutral current sheet that forms the visible **Type I ion tail** (dominated by $\text{CO}^+$ and $\text{H}_2\text{O}^+$ ions).

\subsection{Advection--Diffusion Model of the Type I Plasma Tail}
Downstream in the cometary tail (where flow velocity $u_z$ is roughly constant along the tail axis $z$), the spatial distribution of observed ion column density $N_i(y,z)$ (measured via $\text{CO}^+$ fluorescence photometry) obeys the steady **Shulman--Nazarchuk anisotropic advection--diffusion equation**, derived from the species continuity equation with Fickian cross-field diffusion:
\begin{equation}
    u_z \pderiv{n_i}{z} = D_\perp \frac{\partial^2 n_i}{\partial y^2} + D_\parallel \frac{\partial^2 n_i}{\partial z^2} - \frac{n_i}{\tau_{\text{rec}}},
    \label{eq:shulman_nazarchuk}
\end{equation}
where $D_\perp$ and $D_\parallel$ are the cross-field and field-aligned plasma diffusion coefficients and $\tau_{\text{rec}}$ is the dissociative recombination lifetime. Relating the transverse tail width expansion $\Delta y^2 \sim 4 D_\perp (z / u_z)$ to Bohm or classical collisional diffusion ($D_\perp \propto T_e / B^2$) enables remote photometric estimation of the magnetic field strength $B_{\text{tail}}$ within the cometary plasma tail.
\end{comment}


\subsection{Neutral coma: gas release and ionization}
\subsubsection{General neutral continuity equation}
For a neutral species with number density $n_n$ and velocity $\uvec_n$, the continuity equation is
\begin{equation}
\frac{\partial n_n}{\partial t}+\Div(n_n\uvec_n)=Q_n-L_n,
\label{eq:neutral-cont-general}
\end{equation}
where $Q_n$ and $L_n$ are neutral production and loss rates per unit volume. For a steady, spherically symmetric coma with constant radial expansion speed $u_n$, outside the nucleus and without distributed neutral sources,
\begin{equation}
\frac{1}{r^2}\frac{\dd}{\dd r}\left(r^2 n_n u_n\right)=-\nu_{\rm loss}n_n.
\label{eq:neutral-radial}
\end{equation}
If the total neutral production rate is $Q$ molecules per second, the boundary condition is $4\pi r^2n_nu_n\to Q$ close to the source. The Haser-type solution is
\begin{equation}
\boxed{n_n(r)=\frac{Q}{4\pi u_n r^2}\exp\left(-\frac{r}{L_n}\right)},
\qquad L_n=\frac{u_n}{\nu_{\rm loss}}.
\label{eq:haser}
\end{equation}
Here $L_n$ is the neutral destruction scale length. If losses are negligible over the region of interest, $r\ll L_n$, then
\begin{equation}
\boxed{n_n(r)\simeq\frac{Q}{4\pi u_n r^2}.}
\label{eq:haser-simple}
\end{equation}
A finite nucleus or numerical source region requires a regularized radius, for example $r_{\rm eff}=\sqrt{r^2+r_s^2}$, to avoid the singularity at $r=0$. Such softening is a numerical prescription, not a new physical law.

\subsubsection{Ion production}
For photoionization with frequency $\nu_{\rm ph}$ (in s$^{-1}$), the local ion production rate is
\begin{equation}
q_{\rm ph}=\nu_{\rm ph}n_n.
\label{eq:photoion}
\end{equation}
Including additional ionization channels,
\begin{equation}
q_i=n_n\left(\nu_{\rm ph}+n_e k_{\rm ei}+n_i k_{\rm cx}+\cdots\right),
\label{eq:ionization-full}
\end{equation}
where $k_{\rm ei}$ and $k_{\rm cx}$ denote effective rate coefficients when the corresponding process can be represented in this form. In a simple photoionization-only model,
\begin{equation}
\boxed{q_i=\nu_{\rm ph}n_n,\qquad S_\rho=m_i\nu_{\rm ph}n_n.}
\label{eq:mass-source}
\end{equation}
Combining equations (\ref{eq:haser}) and (\ref{eq:mass-source}) gives
\begin{equation}
S_\rho(r)=\frac{m_iQ\nu_{\rm ph}}{4\pi u_n r^2}\exp\left(-\frac{r}{L_n}\right).
\label{eq:haser-source}
\end{equation}

\subsection{Solar-wind upstream conditions}
\subsubsection{Radial solar-wind approximation}
A first approximation treats the solar wind as a steady radial flow from the Sun:
\begin{equation}
\uvec_{\rm sw}=u_{\rm sw}\,\hat{\boldsymbol r},\qquad
n_{\rm sw}(r)u_{\rm sw}(r)r^2=\mathrm{constant}.
\end{equation}
For nearly constant speed,
\begin{equation}
\boxed{n_{\rm sw}(r)=n_{\rm sw,1AU}\left(\frac{1\,\mathrm{AU}}{r}\right)^2.}
\label{eq:sw-density}
\end{equation}
This is a background scaling; real solar-wind density and speed vary with time and heliographic latitude.

\subsubsection{Parker-spiral magnetic field}
For a steady, radially expanding solar wind and a rotating Sun, a common idealized Parker-spiral field is
\begin{align}
B_r(r)&=B_{r,0}\left(\frac{r_0}{r}\right)^2,\\
B_\phi(r)&=-B_r(r)\frac{\Omega_\odot(r-r_0)}{u_{\rm sw}},
\label{eq:parker-components}
\end{align}
where $\Omega_\odot$ is the solar angular rotation rate and $r_0$ is a reference radius. Far from the reference radius, the magnitude of the azimuthal component can become important. The field magnitude is
\begin{equation}
B=\sqrt{B_r^2+B_\phi^2+B_\theta^2},
\end{equation}
with $B_\theta$ often neglected in the simplest equatorial model. The Parker model is an approximation and should not be interpreted as an in-situ measurement at the comet.

\subsection{Comet-frame solar-wind velocity}
The plasma flow relevant to the interaction is the solar-wind velocity relative to the comet:
\begin{equation}
\boxed{\uvec_{\rm rel}=\uvec_{\rm sw}-\uvec_c,}
\label{eq:relative-velocity}
\end{equation}
where $\uvec_c$ is the comet heliocentric velocity. Its magnitude is
\begin{equation}
u_{\rm rel}=|\uvec_{\rm rel}|.
\end{equation}
A convenient local Cartesian frame takes $\hat{\boldsymbol x}$ along $\uvec_{\rm rel}$, $\hat{\boldsymbol z}$ approximately along the chosen normal to the orbital plane, and $\hat{\boldsymbol y}=\hat{\boldsymbol z}\times\hat{\boldsymbol x}$. The field must be transformed into the same frame before initializing a simulation. A Galilean transformation changes the velocity components but not the magnetic field in non-relativistic MHD to the accuracy normally used here.

\subsection{Full single-fluid ideal-MHD equations}
The following equations describe a conducting plasma when a single-fluid MHD approximation is appropriate. They are written in conservative form and allow explicit source terms.
\subsubsection{Mass conservation}
\begin{equation}
\boxed{\frac{\partial\rhov}{\partial t}+\Div(\rhov\uvec)=S_\rho.}
\label{eq:mhd-mass}
\end{equation}
\subsubsection{Momentum conservation}
\begin{equation}
\boxed{\frac{\partial(\rhov\uvec)}{\partial t}
+\Div\left[\rhov\uvec\uvec+\left(p+\frac{B^2}{2\mu_0}\right)\mathbf I
-\frac{\Bvec\Bvec}{\mu_0}\right]
=\Svec_m.}
\label{eq:mhd-momentum}
\end{equation}
Here $\mathbf I$ is the identity tensor and $\Bvec\Bvec$ is a dyadic product. The source $\Svec_m$ includes momentum carried by newly added mass and any other external forces.
\subsubsection{Total-energy conservation}
Define the total energy density
\begin{equation}
E=\frac{p}{\gamma-1}+\frac{1}{2}\rhov u^2+\frac{B^2}{2\mu_0}.
\label{eq:total-energy}
\end{equation}
Then
\begin{equation}
\boxed{\frac{\partial E}{\partial t}+
\Div\left[\left(E+p+\frac{B^2}{2\mu_0}\right)\uvec
-\frac{(\uvec\cdot\Bvec)\Bvec}{\mu_0}\right]=S_E.}
\label{eq:mhd-energy}
\end{equation}
The source $S_E$ must be consistent with the mass and momentum source model and with the assumed thermal energy of the injected material.
\subsubsection{Induction equation}
For ideal MHD,
\begin{equation}
\boxed{\frac{\partial\Bvec}{\partial t}=\Curl(\uvec\times\Bvec),\qquad \Div\Bvec=0.}
\label{eq:induction-ideal}
\end{equation}
For finite electrical resistivity $\eta$ in the magnetic-diffusivity convention,
\begin{equation}
\frac{\partial\Bvec}{\partial t}=\Curl(\uvec\times\Bvec)-\Curl(\eta\Curl\Bvec).
\label{eq:induction-resistive}
\end{equation}
The ideal equation is the appropriate starting point for an ideal-MHD PLUTO setup; numerical divergence control is still needed to control violations of $\Div\Bvec=0$.
\subsubsection{Equation of state}
For an ideal gas,
\begin{equation}
\boxed{p=(\gamma-1)e_{\rm th},\qquad c_s=\sqrt{\frac{\gamma p}{\rhov}},}
\label{eq:eos}
\end{equation}
where $e_{\rm th}$ is thermal energy density and $c_s$ is the adiabatic sound speed.

\subsection{Mass-loading source terms}
\subsubsection{Conservative source formulation}
Suppose ions are created from neutrals at a rate $q_i$ and enter the plasma with mean injection velocity $\uvec_n$. The mass source is
\begin{equation}
S_\rho=m_iq_i.
\end{equation}
The conservative momentum source associated with the injected mass is
\begin{equation}
\boxed{\Svec_m=S_\rho\uvec_n+\Svec_{\rm ext},}
\label{eq:momentum-source}
\end{equation}
where $\Svec_{\rm ext}$ represents other forces or momentum exchange terms if included. This expression is a fluid-level approximation: pickup-ion acceleration and the associated electromagnetic momentum transfer may require a more detailed multi-fluid or kinetic treatment.

A useful way to understand the velocity response is to expand the conservative equations. Ignoring other forces and using the mass equation,
\begin{equation}
\rhov\left(\frac{\partial\uvec}{\partial t}+\uvec\cdot\grad\uvec\right)
=-\grad p+\Jvec\times\Bvec+S_\rho(\uvec_n-\uvec),
\label{eq:primitive-momentum}
\end{equation}
where $\Jvec=\Curl\Bvec/\mu_0$. The term $S_\rho(\uvec_n-\uvec)$ illustrates mass loading: adding mass with a velocity different from the local plasma velocity changes the bulk flow.

\subsubsection{Energy carried by injected material}
A consistent conservative source may be written schematically as
\begin{equation}
S_E=S_\rho\left(\frac{u_n^2}{2}+\epsilon_{\rm inj}\right)+Q_{\rm heat}-Q_{\rm cool},
\label{eq:energy-source}
\end{equation}
where $\epsilon_{\rm inj}$ is the specific internal energy assigned to the injected material and $Q_{\rm heat}$ and $Q_{\rm cool}$ represent additional heating and cooling rates per unit volume. The exact form depends on the physical closure: ion pickup energy, electron heating, chemical energy, and unresolved kinetic energy should not be counted twice.

For a simplified source model that injects negligible neutral momentum and negligible internal energy in the chosen frame, one may set
\begin{equation}
\Svec_m\simeq\boldsymbol 0,\qquad S_E\simeq 0,
\end{equation}
while retaining $S_\rho>0$. This is a numerical simplification, not a complete physical description of pickup-ion dynamics. Another commonly useful approximation assumes neutrals are nearly at rest in the comet frame, $\uvec_n\simeq\boldsymbol 0$, giving a primitive momentum sink $-S_\rho\uvec$.

\subsubsection{Haser plus photoionization source used in a reduced model}
Combining the simple Haser coma with photoionization gives
\begin{equation}
\boxed{S_\rho(r)=m_i\nu_{\rm ph}\frac{Q}{4\pi u_n r_{\rm eff}^2}
\exp\left(-\frac{r_{\rm eff}}{L_n}\right),\quad
r_{\rm eff}=\sqrt{r^2+r_s^2}.}
\label{eq:reduced-source}
\end{equation}
A two-dimensional Cartesian implementation can use $r=\sqrt{(x-x_c)^2+(y-y_c)^2}$ and apply the source locally to the density, momentum, energy, and tracer equations. The source must be applied consistently with the solver's time integration and operator-splitting strategy.

\subsection{Tracer for cometary material}
A passive mass fraction $C$ can track material originating from the comet. Its conservative equation is
\begin{equation}
\frac{\partial(\rhov C)}{\partial t}+\Div(\rhov C\uvec)=S_{\rho C}.
\label{eq:tracer-general}
\end{equation}
If the injected mass is tagged as cometary material and the upstream solar wind has $C=0$, a natural source is
\begin{equation}
\boxed{S_{\rho C}=S_\rho.}
\label{eq:tracer-source}
\end{equation}
The corresponding non-conservative form, using mass conservation, is
\begin{equation}
\rhov\left(\frac{\partial C}{\partial t}+\uvec\cdot\grad C\right)
=S_{\rho C}-CS_\rho.
\end{equation}
Thus, when $S_{\rho C}=S_\rho$, newly added cometary mass drives the local mass fraction toward unity. The tracer is a diagnostic scalar; it does not by itself alter the MHD dynamics.

\subsection{Non-dimensionalization for numerical simulations}
Choose characteristic scales $L_0$, $V_0$, and $\rhov_0$. Define
\begin{equation}
t_0=\frac{L_0}{V_0},\quad
p_0=\rhov_0V_0^2,\quad
B_0=V_0\sqrt{\mu_0\rhov_0}.
\label{eq:normalization}
\end{equation}
Dimensionless variables are
\begin{equation}
\tilde{\boldsymbol x}=\frac{\boldsymbol x}{L_0},\quad
\tilde t=\frac{t}{t_0},\quad
\tilde\rhov=\frac{\rhov}{\rhov_0},\quad
\tilde{\uvec}=\frac{\uvec}{V_0},\quad
\tilde p=\frac{p}{p_0},\quad
\tilde{\Bvec}=\frac{\Bvec}{B_0}.
\end{equation}
With this choice the dimensionless ideal-MHD equations retain their conservative form and the magnetic coefficient is unity in Alfv\'en units. The dimensionless mass source is
\begin{equation}
\tilde S_\rho=\frac{S_\rho t_0}{\rhov_0}.
\label{eq:dimensionless-source}
\end{equation}
For a Haser-photoionization source,
\begin{equation}
\tilde S_\rho(\tilde r)=
\frac{m_iQ\nu_{\rm ph}t_0}{4\pi u_n\rhov_0L_0^2}
\frac{1}{\tilde r_{\rm eff}^2}
\exp\left(-\frac{L_0\tilde r_{\rm eff}}{L_n}\right).
\end{equation}
Always state the unit system used by the code. In cgs, the magnetic normalization and magnetic-pressure expressions must be written with the corresponding factors of $4\pi$ rather than SI $\mu_0$.



\end{document}